% GENERATED FILE. Edit canonical lessons, then run npm run generate:books.
% canonical-source-hash: fnv1a64:d6258d264cbaa3d2
% canonical-lessons: HI-C394-inam, HI-C394-pratiyogita, HI-C394-pad, HI-C394-qabil, HI-C394-savarna

\chapter{Prizes and Rewards}
\label{ch:hi-prizes-and-rewards}
\hlchaptermodality{394}

\begin{chapteropening}
\textbf{By the end of this chapter:} \emph{I can say what someone will win if they do well.}

Complete the last lesson of chapter 394: I can say what someone will win if they do well.
\end{chapteropening}

\section[\texorpdfstring{\hi{इनाम} (inām)}{इनाम (inām)}]{\hi{इनाम} (inām) — a prize, a reward}

\label{lesson:HI-C394-inam}



\begin{quote}
Before the new one: say the Hindi for a promotion, then the Hindi for a matter of honour.
\end{quote}

\subsection*{You'll want to know: \hi{इनाम}}
\textbf{\hi{इनाम}} — \emph{inām} — \textquotedblleft{}a prize, a reward\textquotedblright{}.

\begin{quote}
\hi{अगर} \hi{तुम} \hi{दौड़} \hi{जीतोगे} \hi{तो} \hi{तुम्हें} \hi{इनाम} \hi{मिलेगा।}
\end{quote}

\begin{quote}
\emph{agar tum dauṛ jītoge to tumhẽ inām milegā.}
\end{quote}

\begin{quote}
If you win the race, you'll get a prize.
\end{quote}

Masculine.

\begin{cousinweb}
\textbf{\hi{इनाम}} (\emph{inām}) is usually traced to Arabic \emph{inʿām}, \textquotedblleft{}a favour, a gift\textquotedblright{}, which reached Hindi through Persian.
\end{cousinweb}

\begin{grammarlens}[title={What you've built}]
One more word for saying what happens if.
\end{grammarlens}

\subsection*{Your turn}
\begin{itemize}
\raggedright
  \item \emph{Say it:} \emph{inām}
  \item \emph{Say it:} \emph{inām}, once more
  \item \emph{Say it:} \emph{agar tum dauṛ jītoge to tumhẽ inām milegā.} — If you win the race, you'll get a prize.
  \item \emph{Recall:} say the Hindi for a promotion, then the Hindi for a matter of honour, then say \emph{inām} again
\end{itemize}

\begin{tcolorbox}[breakable,colback=teal!4,colframe=teal!35!black,title={Before you move on}]
What does \hi{इनाम} mean? (\textquotedblleft{}A prize, a reward\textquotedblright{}.) Say it once more.
\end{tcolorbox}
\section[\texorpdfstring{\hi{प्रतियोगिता} (pratiyogitā)}{प्रतियोगिता (pratiyogitā)}]{\hi{प्रतियोगिता} (pratiyogitā) — a competition}

\label{lesson:HI-C394-pratiyogita}



\begin{quote}
Before the new one: say the Hindi for a matter of honour, then the Hindi for a prize.
\end{quote}

\subsection*{You'll want to know: \hi{प्रतियोगिता}}
\textbf{\hi{प्रतियोगिता}} — \emph{pratiyogitā} — \textquotedblleft{}a competition\textquotedblright{}.

\begin{quote}
\hi{अगर} \hi{कल} \hi{प्रतियोगिता} \hi{होगी} \hi{तो} \hi{पूरा} \hi{स्कूल} \hi{आएगा।}
\end{quote}

\begin{quote}
\emph{agar kal pratiyogitā hogī to pūrā skūl āegā.}
\end{quote}

\begin{quote}
If the competition is tomorrow, the whole school will come.
\end{quote}

Feminine.

\begin{cousinweb}
\textbf{\hi{प्रतियोगिता}} (\emph{pratiyogitā}) is a tatsama, borrowed from Sanskrit \emph{pratiyogitā}, \textquotedblleft{}correlation\textquotedblright{} (Monier-Williams). Monier-Williams does not give the sense of a contest, which is usually taken as a modern Hindi one.
\end{cousinweb}

\begin{grammarlens}[title={What you've built}]
One more word for saying what happens if.
\end{grammarlens}

\subsection*{Your turn}
\begin{itemize}
\raggedright
  \item \emph{Say it:} \emph{pratiyogitā}
  \item \emph{Say it:} \emph{pratiyogitā}, once more
  \item \emph{Say it:} \emph{agar kal pratiyogitā hogī to pūrā skūl āegā.} — If the competition is tomorrow, the whole school will come.
  \item \emph{Recall:} say the Hindi for a matter of honour, then the Hindi for a prize, then say \emph{pratiyogitā} again
\end{itemize}

\begin{tcolorbox}[breakable,colback=teal!4,colframe=teal!35!black,title={Before you move on}]
What does \hi{प्रतियोगिता} mean? (\textquotedblleft{}A competition\textquotedblright{}.) Say it once more.
\end{tcolorbox}
\section[\texorpdfstring{\hi{पद} (pad)}{पद (pad)}]{\hi{पद} (pad) — a post, a position}

\label{lesson:HI-C394-pad}



\begin{quote}
Before the new one: say the Hindi for a prize, then the Hindi for a competition.
\end{quote}

\subsection*{You'll want to know: \hi{पद}}
\textbf{\hi{पद}} — \emph{pad} — \textquotedblleft{}a post, a position\textquotedblright{}.

\begin{quote}
\hi{अगर} \hi{वह} \hi{चुनाव} \hi{जीत} \hi{गया} \hi{तो} \hi{उसे} \hi{मंत्री} \hi{का} \hi{पद} \hi{मिलेगा।}
\end{quote}

\begin{quote}
\emph{agar vah cunāv jīt gayā to use mantrī kā pad milegā.}
\end{quote}

\begin{quote}
If he wins the election, he'll get a minister's post.
\end{quote}

Masculine.

\begin{cousinweb}
\textbf{\hi{पद}} (\emph{pad}) is a tatsama, borrowed from Sanskrit \emph{pada}, \textquotedblleft{}a step, a footing\textquotedblright{} (Monier-Williams).
\end{cousinweb}

\begin{grammarlens}[title={What you've built}]
One more word for saying what happens if.
\end{grammarlens}

\subsection*{Your turn}
\begin{itemize}
\raggedright
  \item \emph{Say it:} \emph{pad}
  \item \emph{Say it:} \emph{pad}, once more
  \item \emph{Say it:} \emph{agar vah cunāv jīt gayā to use mantrī kā pad milegā.} — If he wins the election, he'll get a minister's post.
  \item \emph{Recall:} say the Hindi for a prize, then the Hindi for a competition, then say \emph{pad} again
\end{itemize}

\begin{tcolorbox}[breakable,colback=teal!4,colframe=teal!35!black,title={Before you move on}]
What does \hi{पद} mean? (\textquotedblleft{}A post, a position\textquotedblright{}.) Say it once more.
\end{tcolorbox}
\section[\texorpdfstring{\hi{क़ाबिल} (qābil)}{क़ाबिल (qābil)}]{\hi{क़ाबिल} (qābil) — capable, able}

\label{lesson:HI-C394-qabil}



\begin{quote}
Before the new one: say the Hindi for a competition, then the Hindi for a post.
\end{quote}

\subsection*{You'll want to know: \hi{क़ाबिल}}
\textbf{\hi{क़ाबिल}} — \emph{qābil} — \textquotedblleft{}capable, able\textquotedblright{}.

\begin{quote}
\hi{अगर} \hi{वह} \hi{क़ाबिल} \hi{होता} \hi{तो} \hi{उसे} \hi{यह} \hi{पद} \hi{मिलता।}
\end{quote}

\begin{quote}
\emph{agar vah qābil hotā to use yah pad miltā.}
\end{quote}

\begin{quote}
If he were capable, he would get this post.
\end{quote}

\begin{cousinweb}
\textbf{\hi{क़ाबिल}} (\emph{qābil}) is usually traced to Arabic \emph{qābil}, \textquotedblleft{}receiving, capable\textquotedblright{}, which reached Hindi through Persian.
\end{cousinweb}

\begin{grammarlens}[title={What you've built}]
One more word for saying what happens if.
\end{grammarlens}

\subsection*{Your turn}
\begin{itemize}
\raggedright
  \item \emph{Say it:} \emph{qābil}
  \item \emph{Say it:} \emph{qābil}, once more
  \item \emph{Say it:} \emph{agar vah qābil hotā to use yah pad miltā.} — If he were capable, he would get this post.
  \item \emph{Recall:} say the Hindi for a competition, then the Hindi for a post, then say \emph{qābil} again
\end{itemize}

\begin{tcolorbox}[breakable,colback=teal!4,colframe=teal!35!black,title={Before you move on}]
What does \hi{क़ाबिल} mean? (\textquotedblleft{}Capable, able\textquotedblright{}.) Say it once more.
\end{tcolorbox}
\section[\texorpdfstring{\hi{सँवारना} (sãvārnā)}{सँवारना (sãvārnā)}]{\hi{सँवारना} (sãvārnā) — to put in order, to make (something) look its best}

\label{lesson:HI-C394-savarna}



\begin{quote}
Before the new one: say the Hindi for a post, then the Hindi for capable.
\end{quote}

\subsection*{You'll want to know: \hi{सँवारना}}
\textbf{\hi{सँवारना}} — \emph{sãvārnā} — \textquotedblleft{}to put in order, to make (something) look its best\textquotedblright{}.

\begin{quote}
\hi{अगर} \hi{तुम} \hi{अपना} \hi{कमरा} \hi{सँवारोगे} \hi{तो} \hi{मेहमान} \hi{ख़ुश} \hi{होंगे।}
\end{quote}

\begin{quote}
\emph{agar tum apnā kamrā sãvāroge to mehmān khuś hoṅge.}
\end{quote}

\begin{quote}
If you tidy up your room, the guests will be pleased.
\end{quote}

Past, with \hi{ने} (he or she …): \hi{उसने} \hi{अपना} \hi{घर} \hi{सँवारा}, \emph{usne apnā ghar sãvārā}.

\begin{cousinweb}
\textbf{\hi{सँवारना}} (\emph{sãvārnā}) is a tadbhava, inherited from Sanskrit \emph{saṃvārayati}, \textquotedblleft{}keeps off; puts in order\textquotedblright{} (CDIAL 13021).
\end{cousinweb}

\begin{grammarlens}[title={What you've built}]
One more word for saying what happens if.
\end{grammarlens}

\subsection*{Your turn}
\begin{itemize}
\raggedright
  \item \emph{Say it:} \emph{sãvārnā}
  \item \emph{Say it:} \emph{sãvārnā}, once more
  \item \emph{Say it:} \emph{agar tum apnā kamrā sãvāroge to mehmān khuś hoṅge.} — If you tidy up your room, the guests will be pleased.
  \item \emph{Recall:} say the Hindi for a post, then the Hindi for capable, then say \emph{sãvārnā} again
\end{itemize}

\begin{tcolorbox}[breakable,colback=teal!4,colframe=teal!35!black,title={Before you move on}]
What does \hi{सँवारना} mean? (\textquotedblleft{}To put in order, to make (something) look its best\textquotedblright{}.) Say it once more.
\end{tcolorbox}
